% ============================ AI USE REPORT ============================
\section*{Report on Use of AI Tools}
\addcontentsline{toc}{section}{Report on Use of AI Tools}

\subsection*{Tool used}

Claude Code (Anthropic), version 2.1.322, a command-line agentic coding assistant
based on a large language model, used from 11 to 19 November 2026. The runtime
reported the model identifier \texttt{deepseek-flash[1m]} for the final pass. The tool
operated in a local development environment with file-system access and the ability
to execute Python and typeset \LaTeX{}.

\subsection*{What the tool did}

\begin{enumerate}[leftmargin=1.6em]
\item \textbf{Data engineering.} Wrote and debugged the geohash decoder, the
  cleaning pipeline that converts the 191 MB order file into a typed columnar store,
  and the aggregation to cells, zones and time bands.
\item \textbf{Modelling and computation.} Implemented the gravity model and its
  maximum-likelihood estimator, the user segmentation, the time-space network linear
  program (formulated by the authors, encoded with Gurobi's Python interface), the
  discrete-event simulators, the logit incentive model, the maximal-coverage siting
  heuristic and the Monte-Carlo capacity model.
\item \textbf{Figures.} Produced the five figures included in the paper --- the dataset,
  Q1 validation, the Q2 schedule, Q3 cost against reward and the sensitivity tornado
  --- as vector PDFs with matplotlib; the intermediate files of a superseded figure
  set remain in the repository but are not included in the paper.
\item \textbf{Text.} Drafted the \LaTeX{} manuscript, tables and reference list from
  the computed results under the authors' direction and repeated review, and in the
  final round rewrote the sections that quote the revised results, recompressed the
  layout to the page limit and rewrote the summary.
\item \textbf{Final-round work (18--19 November 2026).} Two sessions of the same tool
  worked in parallel: one drafted written task specifications, acceptance criteria and
  numerical locks; the other wrote and modified the final-round scripts (integer truck
  scheduling, budgeted incentives, capped fences, the same-zone tricycle layer, the
  uncertainty decomposition and the two-platform mismatch experiment), checked each
  delivered number against the stored result files, and rewrote the affected sections,
  the summary and the page layout.
\end{enumerate}

\subsection*{Representative queries and outputs}

The first two examples are from the validation stage and are reproduced from the
working record; their outputs are archived with the code as
\texttt{dest\_validation.json} and \texttt{k\_stability.json}. Examples 3 and 4 are
from the final round and are archived as \texttt{exec\_rows} in the two-tier result
file and as \texttt{cost\_vs\_r.json}.

\noindent\textbf{Example 1: held-out comparison of destination models.}
\begin{quote}
\small
\textbf{Query.} Compare the gravity destination model with distance-only,
attraction-only and uniform benchmarks on held-out days. Fit each model's own
distance-decay parameter; report negative log-likelihood per trip and top-five hit
rate; do not use AIC.

\textbf{Output.} The script fitted five time bands and scored 555{,}675 held-out
trips. Pooled negative log-likelihood was 3.4185 for the full model, 3.4881 for
distance-only, 8.1124 for attraction-only and 8.5685 for uniform; pooled top-five
hit rates were 0.4255, 0.3924, 0.0108 and 0.00095. Because the attraction field is a
smoothed empirical density, the comparison uses out-of-sample loss rather than AIC.
The full model reduced loss by 0.0696 nats (95\% bootstrap interval 0.0685--0.0705)
and raised the top-five hit rate by 3.31 percentage points (3.23--3.38).
\end{quote}

\noindent\textbf{Example 2: stability of the user segmentation.}
\begin{quote}
\small
\textbf{Query.} Test whether the published $k=5$ segmentation is stable. Run
$k=3$--$6$ with at least 20 random initialisations per $k$; report pairwise adjusted
Rand index (mean and minimum), silhouette score, cross-$k$ agreement, and agreement
with the published labels.

\textbf{Output.} The published $k=5$ labels were recovered at ARI 0.994, but across
20 initialisations the pairwise ARI was 0.665 on average and 0.305 at worst. The
best-run silhouette on a 40{,}000-user subsample was 0.256 at $k=5$ and 0.263 at
$k=6$; the mean silhouette on a 10{,}000-user subsample was 0.251 at $k=5$ and 0.246
at $k=6$. The output therefore did not support the claim that the segmentation is
stable across random initialisations.
\end{quote}

\noindent\textbf{Example 3: what the executable tricycle plan actually leaves unserved.}
\begin{quote}
\small
\textbf{Query.} The reported unmet of the two-tier plan is taken from the linear
relaxation; recompute it from the integer trips, adding the shortfall between planned
and delivered bicycle loads cell by cell, and report the layer, the total and the
saving together with the relaxation columns.

\textbf{Output.} The integer trips delivered 1{,}494.80 fewer bicycle loads than the
relaxation; adding that shortfall gives 39{,}984.2332 unserved departures, a layer of
533{,}248.01 and a two-tier total of 617{,}055.75 CNY/day, a saving of
902{,}223.20 (59.385\%) against the no-tricycle baseline, instead of the 59.98\% that
the relaxation supported. The relaxed and integer-executed columns are now reported
side by side.
\end{quote}

\noindent\textbf{Example 4: does the incentive pay inside the tricycle layer?}
\begin{quote}
\small
\textbf{Query.} Re-solve the reroute program inside the per-zone tricycle fixed point
on the reward grid 0--3.00 in 0.05-CNY steps under both source caps; $r=0$ must
reproduce the accepted two-tier plan bit for bit; report the optimum, the net saving,
the payment against the daily budget and the break-even reward.

\textbf{Output.} The $r=0$ run reproduced 975 tricycles, 39{,}984.2332 unserved and
617{,}055.75 CNY/day exactly. The optimum is $r=0.60$ under both source caps:
27{,}680.9 reroutes, a payment of 16{,}608.52 inside the 30{,}000 budget and a net
saving of 18{,}011.46 CNY/day (2.92\%) under the arrivals cap, and 6{,}651.22
(1.08\%) under the surplus cap; the break-even rewards are 1.04 and 1.63. The truck
layer is unchanged by construction.
\end{quote}

\subsection*{What was verified, and how}

\begin{itemize}[leftmargin=1.4em]
\item \textbf{Result numbers are script-generated, with listed exceptions.} Values
  from the preparation and first modelling stages are printed by
  \texttt{95\_numbers.py} from the result files under \texttt{data/}. The final-round
  values are printed from their result files by \texttt{code/t021a\_numbers.py},
  which writes the number lock \texttt{numbers.json} (archived with the final-round
  result files): 724 entries,
  each recording the printed token, the result file and key, a scale factor and a
  caliber, and each verified by re-reading the stored value. The constant 343{,}638
  is still hard-coded in four
  computation scripts and two figure lines: \texttt{20\_q2\_diag.py} line 11,
  used at line 26 for the initial inventory written to \texttt{q2\_state.npz};
  \texttt{24\_cell\_curve.py} line 9, used at line 36 for the cell-level initial
  inventory; \texttt{21\_q2\_opt.py} line 19, where it cancels algebraically;
  \texttt{30\_q3\_incentive.py} line 16, where it is declared but unused; and
  \texttt{90\_figs.py} lines 256 and 258, the cost curve. It is an input to the
  Q2/Q5 runs, not a display value. The Q5 interpretation text in
  \texttt{50\_q5\_platform.py} lines 37--38 is also written by hand. The paper
  keeps 343{,}638 only as the initial-inventory proxy, and does not use it as a
  fleet estimate. The two full-coverage slot bands 600,000--750,000 and
  780,000--900,000 in \texttt{t009\_q4.py} are typed planning bands, used only as
  sanity checks.
\item \textbf{Typed constants of the final round are planning assumptions.} Truck
  scheduling uses 20 km/h, 0.5 CNY per bicycle loaded or unloaded, 600 CNY per
  truck-day and 6 CNY per truck-kilometre, with the fleet scanned over
  $\{0,20,40,60,80\}$ trucks. The tricycle layer uses 20 bicycles per trip, 15 km/h,
  0.5 minutes per bicycle handled, trips of at most 5 km, 300 CNY per tricycle-day,
  0.20 CNY per vehicle-kilometre and one 07:00--20:00 shift. The 150-slot fence cap
  and the 450 m walking radius are typed; the 4,000-site knee is computed by the
  Kneedle rule on the capacity--coverage frontier, not typed. The incentive
  acceptance parameters $\kappa=2$, $\rho=0.35$ and $c_0=0.8$ are assumed (no reward
  was offered in the data), the reward grid is 0.05 CNY, the daily budget is
  30{,}000 CNY, and $c_0$ is re-optimised at 0.8 and 1.2 times its value; a full
  sweep of $\kappa$ and $\rho$ was not run.
\item \textbf{The computation is reproducible in a fixed order, not by a single
  command.} After the preparation scripts 00--04 and \texttt{build\_dist\_matrix.py},
  the numbered stages are run in the order 10 $\to$ 11 $\to$ 12 $\to$ 20 $\to$ 21
  $\to$ \dots $\to$ 95 (\texttt{12\_q1\_od.py} writes the Q1 origin--destination
  surface that the Q2 and Q4 stages read); later stages read the outputs of earlier
  ones (\texttt{q1\_params.json}, \texttt{q2\_state.npz}, and so on). The
  final-round scripts run after that chain, in the order documented with the code.
  \texttt{lib.py} resolves \texttt{data/} and \texttt{figs/} relative to its own
  file, so the tree can be moved. The pipeline is reproducible, but it requires the
  stated order and is not unattended.
\item \textbf{References were checked against Crossref.} Every bibliographic entry
  with a DOI was retrieved from the Crossref API and its authors, year, title,
  journal, volume, issue and pages were compared with the entry. Entries without a
  DOI are books, standards, datasets or news reports cited from their published
  editions or URLs; the McFadden chapter was checked against the standard pagination
  of \emph{Frontiers in Econometrics}.
\item \textbf{External data were fetched from source.} The subway entrances used in
  Q4/Q5 come from a dated Overpass API query (11 November 2026) whose raw response is
  archived with the code, so the 1,097 entrances can be recounted.
\item \textbf{The AI did not invent results.} Where a quantity was not identified by
  the data --- the sampling rate, the initial inventory, the incentive parameters ---
  the manuscript says so and reports a bound or a sweep rather than a point estimate.
\end{itemize}

\subsection*{What the tool did not do}

The research questions, the modelling choices, the interpretation of the results and
the final text were directed and reviewed by the authors. The authors are responsible
for the content of the paper, including any error that survived review.
